% ID: SEL-D0052-OB-E01
% ESTADO: validada. S053.
\textbf{a)} La región queda por encima de $9-3x$ y $(5x-15)/2$, y por debajo de $14-4x/7$, con $x\ge0$. Sobre $x=0$ va de $y=9$ a $y=14$. Las dos cotas inferiores se cruzan en $(3,0)$. La cota inferior ascendente corta a la superior en $(7,10)$. Estos cuatro vértices delimitan el recinto cerrado.
\begin{center}
\setlength{\unitlength}{0.9mm}
\begin{picture}(80,60)(-18,-8)
\qbezier(0,0)(30,0)(60,0)
\qbezier(0,0)(0,23)(0,46)
\put(61,0){$x$}
\put(2,46){$y$}
\qbezier(0,27)(12,13.5)(24,0)
\qbezier(24,0)(40,15)(56,30)
\qbezier(56,30)(28,36)(0,42)
\qbezier(0,42)(0,34.5)(0,27)
\put(-17,26){\scriptsize $(0,9)$}
\put(-17,41){\scriptsize $(0,14)$}
\put(19,-6){\scriptsize $(3,0)$}
\put(57,31){\scriptsize $(7,10)$}
\end{picture}
\end{center}
\textbf{b)} En $(0,9),(0,14),(3,0),(7,10)$, los valores de $F$ son $72,112,42,178$, respectivamente. Como $F$ es lineal y sólo un vértice alcanza el mayor valor, el máximo es $\boxed{178}$ y se alcanza únicamente en $\boxed{(7,10)}$.

\textbf{c)} $(0,10)$ es factible: $10\ge9$, $10\le14$ y $-20\le15$. Da $F=80$, por lo que no es óptimo.
